\section{Senjougahara: sharp affection, trauma, verbal dominance, controlled intimacy}

Сенджогахара уже была примером контролирующей близости (5.1) — здесь добавим то, что делает её случай крайним: степень, в которой её verbal dominance — не стиль общения, а буквально терапевтический инструмент, которым она восстанавливает контроль после травматического периода физической беспомощности.

Вектор: affect — острый, интеллектуально агрессивный; роль — соперница-партнёр, дистанцию которой нужно заслужить; intimacy — контролирующая (5.1); agency — компенсаторно высокая, реакция на период, когда agency у неё не было вообще; conflict — вербальная атака как основной канал, редко физическая.

Деталь: до начала истории она пережила период, когда её тело буквально не подчинялось её воле (сюжетный элемент, объясняющий её позднейшую одержимость контролем), и её острый язык — прямая компенсация этого опыта, а не врождённая черта.

Для книги это ценный контрпример к главе 4: там мы видели баги, которые балансируют силу (Мисато, 4.7), а Сенджогахара показывает обратный случай — силу, которая сама является компенсацией за пережитую беспомощность. Для Рин это не прямая модель (у неё нет истории физической беспомощности), но структура полезна: любая острота должна иметь конкретный источник, а не быть просто "стилем".
